% !TEX root = ../main.tex

\chapter*{Conclusion}
\addcontentsline{toc}{chapter}{Conclusion}

Ce stage au sein d'ISI.nc, réalisé pour le compte du client FPG, m'a
permis de mettre en pratique l'ensemble de mes compétences en
développement full stack sur un projet réel, complexe et exigeant. La
mission qui m'a été confiée  -  la conception et le développement d'un
système d'historisation des modifications  -  s'est révélée
particulièrement formatrice, tant sur le plan technique (architecture
backend, conception de composants frontend réutilisables, débogage
transverse) que sur le plan méthodologique (Scrum, collaboration avec une
équipe produit externe).

Au regard des objectifs pédagogiques fixés en début de stage, je considère
que les compétences visées ont été atteintes, voire dépassées sur
certains aspects techniques que je ne maîtrisais pas initialement, comme
l'historisation avancée avec \texttt{django-simple-history}.

Cette expérience conforte mon projet professionnel de poursuivre dans le
développement full stack, et m'a permis de mesurer concrètement
l'importance de la rigueur, de la communication et de la démarche de
débogage méthodique dans un contexte professionnel.

Sur le plan des perspectives, la fonctionnalité d'historisation
développée pourrait être enrichie par des fonctionnalités de recherche
et de filtrage avancées dans l'historique, ou encore par des exports
dédiés à des fins d'audit, ouvrant la voie à de futures évolutions pour
l'équipe qui poursuivra ce projet.

\newpage
